\documentclass[12pt,a4paper]{article}  % Déclare la classe du document.

\newcommand{\chaptermark}[2]{\og #2 \fg{} (#1)} % ligne uniquement pour éviter une erreur en cas d'article
\usepackage{mathrsfs}
\include{../1ere-preambule}

\usepackage{epic}   % extension et simplification de  l'environnement picture
\usepackage{graphicx}    % permet d'inclure des graphique externe dans un document
\usepackage{arydshln}   %permet d'avoir des lignes en pointillés

\newcommand{\lignes}[1]{
\setlength{\unitlength}{1mm}
\vskip 0,4cm\noindent
\multido{\i=0+1}{#1}{%
     \noindent
     \begin{picture}(150,5)
     \linethickness{0,005mm}
     \put(-5,0){\dashline{1}(175,5)(5,5)}
     \end{picture}\vskip 0,001cm
     }
\vskip 0cm
}

\rhead{\textcolor{red}{Correction - DS n$^{\circ}$06 - Sujet A - 15/05/2022}}

\begin{document}
\graphicspath{{../images/}}

\begin{center}
\fbox{
\begin{minipage}[]{15cm}
\begin{center}
\LARGE{\rule{0.5cm}{0pt}\rule[-0.25cm]{0pt}{0.9cm}
\textcolor{red}{Correction - DS n$^{\circ}$06 - Sujet A - Fonction exponentielle\\
Application du produit scalaire}
}
\end{center}
\end{minipage}
}

\end{center}


%\begin{center}\textbf{1h - Calculatrice autorisée - Barème indicatif}\end{center}
%
%\begin{center}\textbf{Les élèves avec un tier-temps ne traitent pas les questions avec  le symbole \ding{46}}\end{center}
%
%\begin{center}
%%On veillera à ENCADRER ses résultats et soigner sa copie.\\
%Les résultats doivent être justifiés par des calculs (au moins 1 étape intermédiaire).
%\end{center}
%
%\begin{center}\textbf{Les élèves avec un tier-temps ne traitent pas les questions avec  le symbole \ding{46}}\end{center}

\exerciceDS{ - Au golf - ( 2 points )}
\begin{minipage}[c]{.6\linewidth}
Sur un green de golf, trois positions de la balle sont repérées en $\mathrm{M},$ N et P sur le schéma ci-contre. Le triangle MNP est tel que MN=8, $\mathrm{MP}=5$ et $\widehat{M}=53^{\circ}$.\\
Calculer la valeur exacte de NP, puis son arrondi à 0,1 près.
\end{minipage}
\hfill
\begin{minipage}[c]{.35\linewidth}
\begin{center}
   \includegraphics[scale=0.8]{1.05.DS-IMG-01} 
\end{center}

\end{minipage}

\begin{corrBox}
	D'après le théorème d'Al-Kashi et avec les notations de la figure, on a :\\
	$NP^2=NM^2+MP^2-2\times NM\times MP \times \cos(\widehat{NMP})$\\
	$NP^2=64+25-80\cos(53^{\circ})$ d'où\\
	$NP=\sqrt{89-80\cos(53^{\circ})}\approx6,4$
\end{corrBox}

\exerciceDS{ - Angles dans un triangle - ( 3 points )}
Dans un repère orthonormé du plan, on considère les points $\mathrm{A}(-2 ;3), \mathrm{B}(1 ; 4)$ et $\mathrm{C}(4 ; 1)$.\\
Déterminer les angles du triangle $A B C$ à 0,1 degré près.\\

\begin{corrBox}
	On utilise le théorème d'AL-Kashi : Dans un triangle $\mathrm{ABC},$ on a, avec les notations de la figure :
	$$
	a^{2}=b^{2}+c^{2}-2 b c \cos \hat{A} \iff \hat{A}=\cos^{-1}\left(\dfrac{b^2+c^2-a^2}{2bc}\right)
	$$
	\begin{center}
		\includegraphics[scale=0.7]{1.04.IMG-03}
	\end{center}
	On commence donc par calculer les 3 longueurs:\\
	$\Coord{AB}{x_B-x_A}{y_B-y_A}=\Coord{AB}{1-(-2)}{4-3}=\Coord{AB}{3}{1} \Longrightarrow AB=\sqrt{3^2+1^2}=\sqrt{10}$\\
	$\Coord{AC}{x_C-x_A}{y_C-y_A}=\Coord{AC}{4-(-2)}{1-3}=\Coord{AC}{6}{-2} \Longrightarrow AC=\sqrt{6^2+(-2)^2}=\sqrt{40}=2\sqrt{10}$\\
	$\Coord{BC}{x_C-x_B}{y_C-y_B}=\Coord{BC}{4-1}{1-4}=\Coord{BC}{3}{-3} \Longrightarrow BC=\sqrt{3^2+(-3)^2}=\sqrt{18}=3\sqrt{2}$\\
	On peut maintenant calculer les angles:\\
	$\widehat{BAC}=\cos^{-1}\left(\dfrac{AC^2+AB^2-BC^2}{2\times AC\times AB}\right)=\cos^{-1}\left(\dfrac{40+10-18}{2\times \sqrt{40}\times \sqrt{10}}\right)\approx36,9^{^\circ}$\\
	$\widehat{ABC}=\cos^{-1}\left(\dfrac{BC^2+AB^2-AC^2}{2\times BC\times AB}\right)=\cos^{-1}\left(\dfrac{18+10-40}{2\times \sqrt{18}\times \sqrt{10}}\right)\approx116.6^{^\circ}$\\
	$\widehat{ACB}=180-\widehat{BAC}-\widehat{ABC}=180-36.9-116.6=26.5^{^\circ}$
\end{corrBox}


\exerciceDS{ - Médiatrice dans un triangle - ( 8 points )}
On considère le triangle $\mathrm{ABC}$ tel que $\mathrm{A}(-2 ;3), \mathrm{B}(1 ; 4)$ et $\mathrm{C}(4 ; 1)$.
\begin{enumMet}
    \item Montrer que les coordonnées du milieu M du segment $[BC]$ sont $M(2,5;2,5)$\\
    \begin{corrBox}
    	$x_M=\dfrac{x_B+x_C}{2}=\dfrac{1+4}{2}=\dfrac{5}{2}=2,5\quad \quad \quad \quad y_M=\dfrac{y_B+y_C}{2}=\dfrac{4+1}{2}=\dfrac{5}{2}=2,5$
    \end{corrBox}
    \item Déterminer une équation de la médiatrice  du segment $[BC]$ dans le triangle $\mathrm{ABC}$ (la médiatrice est la perpendiculaire à $[BC]$ passant par M).\\
    \begin{corrBox}
    	La médiatrice du segment $[BC]$ est la droite passant par $M$ et perpendiculaire à $(BC)$. Le vecteur $\vv{BC}$ est donc un vecteur normal à la médiatrice du segment $[BC]$.\\
    	Comme $\Coord{BC}{4-1}{1-4}=\Coord{BC}{3}{-3}$, une équation de la hauteur est de la forme $3x-3y+c=0$.\\
    	Or $M$ appartient à la médiatrice, donc $3x_{M}-3y_{M}+c=0$, donc $3\times 2,5-3\times2,5+c=0$, donc $7,5-7,5+c=0$, donc $c=0$.\\
    	Une équation de la médiatrice du segment $[BC]$ est alors $3x-3y=0$, ce qui équivaut à $x-y=0$.
    \end{corrBox}
    \item \ding{46} Montrer que les coordonnées du milieu N du segment $[AC]$ sont $N(1;2)$\\
    \begin{corrBox}
    	$x_N=\dfrac{x_A+x_C}{2}=\dfrac{-2+4}{2}=\dfrac{2}{2}=1\quad \quad \quad \quad y_N=\dfrac{y_A+y_C}{2}=\dfrac{3+1}{2}=\dfrac{4}{2}=2$
    \end{corrBox}
    \item Déterminer une équation de la médiatrice  du segment $[AC]$ dans le triangle $\mathrm{ABC}$.\\
    \begin{corrBox}
    	La médiatrice du segment $[AC]$ est la droite passant par $N$ et perpendiculaire à $(A C)$. Le vecteur $\vv{A C}$ est donc un vecteur normal à la médiatrice du segment $[AC]$.\\
    	Comme $\Coord{AC}{4-(-2)}{1-3}=\Coord{AC}{6}{-2}$, une équation de la hauteur est de la forme $6x-2y+c=0$.\\
    	Or $N$ appartient à la médiatrice, donc $6x_{N}-2y_{N}+c=0$, donc $6\times 1-2\times2+c=0$, donc $6-4+c=0$, donc $c=-2$.\\
    	Une équation de la médiatrice du segment $[AC]$ est alors $6x-2y-2=0$, ce qui équivaut à $3x-y-1=0$.
    \end{corrBox}
    \item Déterminer les coordonnées du point d'intersection $\mathrm{O}$ des deux médiatrices.\\
    \begin{corrBox}
    	Les 2 médiatrices ont pour équations réduites $y=x$ et $y=3x-1$. On égalise les 2 équations:\\
    	$x=3x-1\iff x-3x=-1\iff -2x=-1\iff x=\dfrac{1}{2}=0,5$\\
    	On remplace $x$ dans la 1ère équation et on trouve $y=0,5$.\\
    	Le point O a pour coordonnées $O(0,5;0,5)$
    \end{corrBox}
    \item \ding{46} Montrer que les coordonnées du milieu P du segment $[AB]$ sont $P(-0,5;3,5)$\\
    \begin{corrBox}
    	$x_P=\dfrac{x_A+x_B}{2}=\dfrac{-2+1}{2}=\dfrac{-1}{2}=-0,5\quad \quad \quad \quad y_P=\dfrac{y_A+y_B}{2}=\dfrac{3+4}{2}=\dfrac{7}{2}=3,5$
    \end{corrBox}
    \item Calculer le produit scalaire $\vv{\mathrm{AB}} \cdot \vv{\mathrm{OP}}$. Quel résultat connu retrouve-t-on ?\\
    \begin{corrBox}
    	$\Coord{OP}{-0,5-0,5}{3,5-0,5}$, donc $\Coord{OP}{-1}{3}$ et $\Coord{AB}{1-(-2)}{4-3}$, donc $\Coord{AB}{3}{1}$.\\
    	Alors $\vv{A B} \cdot \vv{OP}=3 \times (-1)+1 \times3=-3+3=0$.\\
    	Les vecteurs $\vv{OP}$ et $\vv{AB}$ sont orthogonaux, donc $(OP)$ est la médiatrice du segment $[AB]$. On en déduit que les médiatrices sont concourantes en $O$.
    \end{corrBox}
\end{enumMet}


\exerciceDS{ - Équations cartésiennes - ( 3 points )}
Dans chaque cas, déterminer une équation cartésienne de la droite $d$ passant par le point $\mathrm{A}$ et perpendiculaire à la droite (BC).
\begin{enumMet}
    \item \ding{46} $\mathrm{A}(5 ; 4)$ et $\Coord{u}{-2}{4}$ un vecteur directeur de $(\mathrm{BC})$.\\
    \begin{corrBox}
    	Comme $\Coord{u}{-2}{4}$ est un vecteur directeur de la droite $(B C)$, il est un vecteur normal à la droite $d$.\\
    	Une équation cartésienne est de la forme $-2x+4 y+c=0$.\\
    	Or, le point $A$ appartient à la droite $d$, donc $-2x_{A}+4y_{A}+c=0$, donc $-2 \times 5+4 \times 4+c=0$, donc $c=-6$.\\
    	Une équation de la droite $d$ est donc $-2 x+4 y-6=0$ ou encore $-x+2y-3=0$.
    \end{corrBox}
    \item $\mathrm{A}(0 ; 2), \mathrm{B}(3 ;-1)$ et $\mathrm{C}(-4 ; 2)$.\\
    \begin{corrBox}
    	Comme $\Coord{BC}{-7}{3}$  est un vecteur directeur de la droite $(B C)$, il est un vecteur normal à la droite $d$.\\
    	Une équation cartésienne est de la forme $-7 x+3y+c=0$.\\
    	Or, le point $A$ appartient à la droite $d$, donc $-7 x_{A}+3y_{A}+c=0$, donc $-7 \times 0+3 \times 2+c=0$, donc $c=-6$.\\
    	Une équation de la droite $d$ est donc $-7 x+3 y-6=0$        
    \end{corrBox}
    \item $\mathrm{A}(-2 ; 3)$ et la droite $(\mathrm{BC})$ est parallèle à l'axe des abscisses.\\
    \begin{corrBox}
    	La droite $(B C)$ étant horizontale, un vecteur directeur de cette droite et un vecteur normal à $d$ est $\Coord{u}{1}{0}$.\\
    	Une équation cartésienne de $d$ est de la forme $x+c=0$.\\
    	Or, le point $A$ appartient à la droite $d$, donc $x_{A}+c=0$, donc $-2+c=0$, donc $c=2$.\\
    	Une équation de la droite $d$ est donc $x+2=0$.        
    \end{corrBox}
\end{enumMet}

\exerciceDS{ - simplification - ( 4 points )}
$x$ est un réel quelconque. Simplifier au maximum les expressions suivantes.
%\begin{multicols}{2}
\begin{enumMet}
	\item $a(x)=\mathrm{e}^{2 x} \times\left(\mathrm{e}^x\right)^2 \times \mathrm{e}^{-3 x}$ \quad \corr{$a(x)=\mathrm{e}^{2 x} \times\mathrm{e}^{2 x} \times \mathrm{e}^{-3 x}=\mathrm{e}^{2x+2x-3x}=\mathrm{e}^{x}$} 
	\item $b(x)=\dfrac{\mathrm{e}^{x^2}}{\mathrm{e}^x}$\quad \corr{$b(x)=\mathrm{e}^{x^2-x}$}
	\item $c(x)=\dfrac{\mathrm{e}^{x-1} \times \mathrm{e}^{4 x}}{\mathrm{e}^x}$\quad \corr{$c(x)=\mathrm{e}^{x-1+4x-x} =\mathrm{e}^{4x-1}$}
	\item \ding{46} $d(x)=\dfrac{\mathrm{e}^{-2 x}}{\mathrm{e}^{-3 x} \times \mathrm{e}^{x+1}}$\quad \corr{$d(x)=\mathrm{e}^{-2 x-(-3x)-(x+1)}=\mathrm{e}^{-1}$}
\end{enumMet}
%\end{multicols}

\exerciceDS{ - dérivation - ( 4 points )}
Calculer la dérivée de la fonction $f$ définie sur $\mathbb{R}$.
\begin{multicols}{2}
	\begin{enumMet}
		\item $f(x)=4 \mathrm{e}^x-8 x^2+2$ \quad \corr{$f'(x)=4 \mathrm{e}^x-16x$}
		\item \ding{46} $f(x)=2x-3 \mathrm{e}^x+1$ \quad \corr{$f'(x)=2-3 \mathrm{e}^x$}
		\item $f(x)=2\mathrm{e}^x+\mathrm{e}^3$ \quad \corr{$f'(x)=2 \mathrm{e}^x$}
		\item $f(x)=2x \mathrm{e}^x$ \quad \corr{$f'(x)=2\mathrm{e}^x+2x\mathrm{e}^x=(2+2x)\mathrm{e}^x$}
	\end{enumMet}
\end{multicols}

\exerciceDS{ - Résolution d'équations et d'inéquations - ( 8 points )}
Résoudre dans \R.
%\begin{multicols}{2}
	\begin{enumMet}
		\item $\mathrm{e}^{2x-3}=\mathrm{e}$ \quad \corr{$\iff 2x-3=1\iff x=2 \iff S=\{2\}$}
		\item $\mathrm{e}^{-2x^2}=\dfrac{1}{\mathrm{e}^{-8}}$ \quad \corr{$\iff\mathrm{e}^{-2x^2}=\mathrm{e}^8 \iff -2x^2=8\iff S=\varnothing$ car $x^2>0$}
		\item \ding{46} $\mathrm{e}^{2x^2-6x}=1$\quad \corr{$\iff 2x^2-6x=0\iff 2x(x-3)=0 \iff S=\{0;3\}$}
		\item $5+\mathrm{e}^x=1$ \quad \corr{$\iff \mathrm{e}^x=-4\iff S=\varnothing$ car $\mathrm{e}^x>0$ pour tout $x\in\R$}
		\item $\mathrm{e}^{x+1}<1$ \quad \corr{$\iff x+1<0\iff x<-1 \iff S=]-\infty;-1[$}
		\item $\mathrm{e}^{-2 x+5} \geqslant 0$ \quad \corr{$\iff S=\R$ car $\mathrm{e}^x>0$ pour tout $x\in\R$}
		\item $-3 \mathrm{e}^{x^2-4}>4$ \quad \corr{$\iff \mathrm{e}^{x^2-4}<-\dfrac{4}{3}\iff S=\varnothing$ car $\mathrm{e}^x>0$ pour tout $x\in\R$}
		\item $\mathrm{e}^{x+4} \leqslant \dfrac{1}{\mathrm{e}^{3 x}}$\\
		\corr{$\iff \mathrm{e}^{x+4} \leqslant \mathrm{e}^{-3 x}\iff x+4\leqslant-3x\iff x+3x\leqslant-4 \iff x\leqslant-1\iff S=]-\infty;-1[$}
	\end{enumMet}
%\end{multicols}



\end{document}